\documentclass[11pt, a4paper]{article}

\usepackage[utf8]{inputenc}
\usepackage[T1]{fontenc}
\usepackage[french]{babel}

% Math packages (amsmath is required for the 'align' environment)
\usepackage{amsmath}
\usepackage{amssymb}
\usepackage{amsfonts}
\usepackage{mathtools}

% Layout and typography
\usepackage[margin=2.5cm]{geometry}
\usepackage{microtype}
\usepackage{hyperref, csquotes}

\begin{document}

  \noindent Longueur des tiges: $l$\\
  longueur: $X=xl$\\
  largeur: $Y=yl$\\
  hauteur: $Z=zl$


  \begin{align}
    nbr_{\vec{x}} &= x\times y \times z\label{eq:1}\\
    nbr_{\vec{y}} &= x\times y \times z\label{eq:2}\\
    nbr_{\vec{z}} &= x\times y \times z\label{eq:3}\\
    nrb &= nbr_{\vec{x}} + nbr_{\vec{y}} + nbr_{\vec{z}} = 3xyz
  \end{align}
  \eqref{eq:1}, \eqref{eq:2} et \eqref{eq:3} sont egaux par symmetrie. Deplus ils n'y a pas de colisions car chaque barre est sur un axe différent.

  donc
  \begin{equation}
    nbr(l, X, Y, Z) \approx 3\frac{XYZ}{l^3}
  \end{equation}
  avec une égalité lorsque $X$, $Y$ et $Z$ sont \enquote{divisibles} par $l$.

  Aussi la longueur total de tube nécessaire pour un meuble est:
  \begin{equation}
    \mathcal{L}\left(l, X, Y, Z\right)=nbr(l, X, Y, Z)\times l \approx 3\frac{XYZ}{l^2}
  \end{equation}

  Soit $r_i$ le rayon interieur des tubes et $r_e$ le rayon exterieur. Le volume d'un tube $v$ est
  \begin{equation}
    v=\left( r_e^2\pi - r_i^2\pi \right)l = \pi l\left( r_e^2 - r_i^2 \right)
  \end{equation}

  Soit $\rho$ la masse volumique de l'acier inoxidable. Un tube pèse $v\rho$. Donc le poid total du meuble est
  \begin{equation}
    pds = 3\frac{XYZ}{l^3}v\rho=3\frac{\pi\rho}{l^2}\left( r_e^2 - r_i^2 \right)XYZ
  \end{equation}

  Soit $L$ la longueur des tubes fournis et $\mathcal{L}_L$ le nombre de ces tubes:
  \begin{equation}
    \mathcal{L}_L = \frac{nbr\left( l, X, Y, Z \right)}{\left\lfloor \frac{L}{l} \right\rfloor}
  \end{equation}

  Le volume d'un long tube est
  \begin{equation}
    v_L=\pi L\left( r_e^2 - r_i^2 \right)
  \end{equation}

  Ce qui nous done un poid de commande de:
  \begin{equation}
    pds_L=\mathcal{L}_Lv_L\rho=3\frac{\pi\rho}{l^3\left\lfloor \frac{L}{l} \right\rfloor}\left( r_e^2 - r_i^2 \right)XYZ
  \end{equation}

  ensuite
  \begin{equation}
    pds_L \geq 3\frac{\pi\rho}{l^3\left(\frac{L}{l} + 1 \right)}\left( r_e^2 - r_i^2 \right)XYZ=3\frac{\pi\rho}{l^2\left(L - l \right)}\left( r_e^2 - r_i^2 \right)XYZ
  \end{equation}
  donc (puisque tout est positif)
  \begin{equation}
    l^2\left(L + l \right) \geq 3\frac{\pi\rho}{pds_L}\left( r_e^2 - r_i^2 \right)XYZ
  \end{equation}

  Selon Wolfram alpha
  \begin{verbatim}
    x^2 (80 cm (centimeters) - x)>=
      3 pi*(8030 kg/m^3 (kilograms per cubic meter))/(10 kg (kilogram))
         ((7 mm (millimeters))^2 - (5.8 mm (millimeters))^2)
         *70 cm (centimeters)
         *70 cm (centimeters)
         *70 cm (centimeters)
  \end{verbatim}
  est 
  \begin{verbatim}
    x^2 (80 cm (centimeters) - x)>=0.04 m^2 (square meters)
  \end{verbatim}
  Ce qui nous done (en mètres)
  \begin{verbatim}
    0.276393<=x<=0.723607
  \end{verbatim}

\end{document}